\documentclass[12pt]{article}

% Match the earlier reports using locally available fonts.
\usepackage[T1]{fontenc}
\usepackage{times}
\usepackage{microtype}
\usepackage{amsmath}
\usepackage{amsfonts}
\usepackage[parfill]{parskip}
\setlength{\parindent}{0pt}
\setlength{\parskip}{2ex plus 0.2ex minus 0.2ex}
\usepackage[letterpaper,margin=1in]{geometry}
\usepackage[hidelinks]{hyperref}

\begin{document}

\begin{flushleft}
\textbf{Week 8: Exponential Families, Conjugate Priors, and GLMs} \\
Jaisel Singh (js6897) \\
November 4, 2026
\end{flushleft}

\vspace{0.1in}

What stood out to me in Bishop and Efron was how much familiar distributions
share once they are written in exponential-family form. For a regular
one-parameter family, we can write
\mbox{$p(y\mid\eta)=h(y)\exp\{\eta T(y)-A(\eta)\}$}.
The sufficient statistic summarizes the information in an observation about
$\eta$. The log normalizer also carries information about uncertainty:
$A'(\eta)$ gives the mean of $T(Y)$, and $A''(\eta)$ gives its variance.
I found this more useful than
memorizing separate formulas for each distribution. The same structure
organizes how observations accumulate information and how uncertainty changes
with the parameter.

Diaconis and Ylvisaker made conjugacy more interesting to me than a table of
matched distributions. Under regularity conditions, they characterize familiar
conjugate priors through linear posterior expectations of mean parameters.
The beta--Bernoulli example makes this concrete: successes update the first
beta parameter and failures update the second, so the posterior mean is a
weighted average of the prior mean and the observed success rate. Their
introduction also shows that closure under updating alone does not uniquely
select the usual conjugate family. I found that a useful reminder that
conjugacy is a computational convenience, and choosing its hyperparameters
still expresses substantive prior beliefs.

The GLM framework connects these ideas back to week 2. We specify an
exponential-family response distribution and a link
\mbox{$g(\mu_i)=x_i^{\mathsf{T}}\beta$}, where
$\mu_i=\mathbb{E}[Y_i\mid x_i]$. It is the transformed conditional mean that
is linear in the coefficients; the response itself can be binary or a count.
The logit link for Bernoulli outcomes and the log link for Poisson outcomes
make this distinction clear. However, an exponential-family likelihood does
not automatically give a convenient conjugate posterior over regression
coefficients: a Gaussian prior on logistic-regression weights still requires
additional inference machinery. Efron's discussion of overdispersion also
made me think about the cost of simplicity. My remaining question is how to
decide when a convenient mean--variance relationship has become too restrictive
to give trustworthy uncertainty estimates.

\vspace{0.05in}
{\footnotesize
\textit{Readings:} Bishop (2006), \emph{Pattern Recognition and Machine Learning}, Section 2.4;
Diaconis and Ylvisaker (1979), \href{https://www.cs.columbia.edu/~blei/fogm/readings/DiaconisYlvisaker1979.pdf}{\emph{Conjugate Priors for Exponential Families}};
Efron (2018), \href{https://www.cs.columbia.edu/~blei/fogm/readings/Efron2018.pdf}{\emph{Exponential Families in Theory and Practice}}.
\par}

\end{document}

%%% Local Variables:
%%% mode: latex
%%% TeX-master: t
%%% End:
